% GENERATED FILE. Edit canonical lessons, then run npm run generate:books.
% canonical-source-hash: fnv1a64:77ccc191d0016daf
% canonical-lessons: MR-C422-mandai, MR-C422-gongat, MR-C422-kolahal, MR-C422-gajbajlela, MR-C422-vardal

\chapter{Noise in the Market}
\label{ch:mr-noise-in-the-market}
\hlchaptermodality{422}

\begin{chapteropening}
\textbf{By the end of this chapter:} \emph{I can talk about a vegetable market, noise, an uproar and a bustling street.}

Complete the last lesson of chapter 422: I can talk about a vegetable market, noise, an uproar and a bustling street.
\end{chapteropening}

\section[\texorpdfstring{\mr{मंडई} (maṇḍaī)}{मंडई (maṇḍaī)}]{\mr{मंडई} (maṇḍaī) — a vegetable market}

\label{lesson:MR-C422-mandai}



\begin{quote}
Before the new one: say the Marathi for a village gate, then the Marathi for a boundary.
\end{quote}

\subsection*{You'll want to know: \mr{मंडई}}
\textbf{\mr{मंडई}} — \emph{maṇḍaī} — \textquotedblleft{}a vegetable market\textquotedblright{}.

\begin{quote}
\mr{आई} \mr{रोज} \mr{सकाळी} \mr{मंडईत} \mr{जाते}.
\end{quote}

\begin{quote}
\emph{āī roj sakāḷī maṇḍaīt jāte.}
\end{quote}

\begin{quote}
Mother goes to the vegetable market every morning.
\end{quote}

Feminine.

A covered market for vegetables and fruit, as in Pune.

\begin{cousinweb}
\textbf{\mr{मंडई}} (\emph{maṇḍaī}) is a Marathi word for a covered vegetable market; its origin is not settled.
\end{cousinweb}

\begin{grammarlens}[title={What you've built}]
One more word for getting around.
\end{grammarlens}

\subsection*{Your turn}
\begin{itemize}
\raggedright
  \item \emph{Say it:} \emph{maṇḍaī}
  \item \emph{Say it:} \emph{maṇḍaī}, once more
  \item \emph{Say it:} \emph{āī roj sakāḷī maṇḍaīt jāte.} — Mother goes to the vegetable market every morning.
  \item \emph{Recall:} say the Marathi for a village gate, then the Marathi for a boundary, then say \emph{maṇḍaī} again
\end{itemize}

\begin{tcolorbox}[breakable,colback=teal!4,colframe=teal!35!black,title={Before you move on}]
What does \mr{मंडई} mean? (\textquotedblleft{}A vegetable market\textquotedblright{}.) Say it once more.
\end{tcolorbox}
\section[\texorpdfstring{\mr{गोंगाट} (goṅgāṭ)}{गोंगाट (goṅgāṭ)}]{\mr{गोंगाट} (goṅgāṭ) — noise; a din}

\label{lesson:MR-C422-gongat}



\begin{quote}
Before the new one: say the Marathi for a boundary, then the Marathi for a vegetable market.
\end{quote}

\subsection*{You'll want to know: \mr{गोंगाट}}
\textbf{\mr{गोंगाट}} — \emph{goṅgāṭ} — \textquotedblleft{}noise; a din\textquotedblright{}.

\begin{quote}
\mr{वर्गात} \mr{खूप} \mr{गोंगाट} \mr{होता}.
\end{quote}

\begin{quote}
\emph{vargāt khūp goṅgāṭ hotā.}
\end{quote}

\begin{quote}
There was a lot of noise in the classroom.
\end{quote}

Masculine.

\mr{गोंगाट} \mr{करणे}, to make a din.

\begin{cousinweb}
\textbf{\mr{गोंगाट}} (\emph{goṅgāṭ}) is a Marathi word that imitates loud chatter; its origin is not settled.
\end{cousinweb}

\begin{grammarlens}[title={What you've built}]
One more word for getting around.
\end{grammarlens}

\subsection*{Your turn}
\begin{itemize}
\raggedright
  \item \emph{Say it:} \emph{goṅgāṭ}
  \item \emph{Say it:} \emph{goṅgāṭ}, once more
  \item \emph{Say it:} \emph{vargāt khūp goṅgāṭ hotā.} — There was a lot of noise in the classroom.
  \item \emph{Recall:} say the Marathi for a boundary, then the Marathi for a vegetable market, then say \emph{goṅgāṭ} again
\end{itemize}

\begin{tcolorbox}[breakable,colback=teal!4,colframe=teal!35!black,title={Before you move on}]
What does \mr{गोंगाट} mean? (\textquotedblleft{}Noise; a din\textquotedblright{}.) Say it once more.
\end{tcolorbox}
\section[\texorpdfstring{\mr{कोलाहल} (kolāhal)}{कोलाहल (kolāhal)}]{\mr{कोलाहल} (kolāhal) — an uproar}

\label{lesson:MR-C422-kolahal}



\begin{quote}
Before the new one: say the Marathi for a vegetable market, then the Marathi for noise.
\end{quote}

\subsection*{You'll want to know: \mr{कोलाहल}}
\textbf{\mr{कोलाहल}} — \emph{kolāhal} — \textquotedblleft{}an uproar\textquotedblright{}.

\begin{quote}
\mr{बातमी} \mr{ऐकून} \mr{सभेत} \mr{कोलाहल} \mr{माजला}.
\end{quote}

\begin{quote}
\emph{bātmī aikūn sabhet kolāhal mājlā.}
\end{quote}

\begin{quote}
Hearing the news, an uproar broke out in the meeting.
\end{quote}

Masculine.

A literary word; \mr{कोलाहल} \mr{माजणे}, an uproar breaks out.

\begin{cousinweb}
\textbf{\mr{कोलाहल}} (\emph{kolāhal}) is a tatsama, taken from Sanskrit \emph{kolāhala}, \textquotedblleft{}a great din\textquotedblright{} (CDIAL 3534), where CDIAL also lists the inherited Marathi \emph{kolhāḷ}.
\end{cousinweb}

\begin{grammarlens}[title={What you've built}]
One more word for getting around.
\end{grammarlens}

\subsection*{Your turn}
\begin{itemize}
\raggedright
  \item \emph{Say it:} \emph{kolāhal}
  \item \emph{Say it:} \emph{kolāhal}, once more
  \item \emph{Say it:} \emph{bātmī aikūn sabhet kolāhal mājlā.} — Hearing the news, an uproar broke out in the meeting.
  \item \emph{Recall:} say the Marathi for a vegetable market, then the Marathi for noise, then say \emph{kolāhal} again
\end{itemize}

\begin{tcolorbox}[breakable,colback=teal!4,colframe=teal!35!black,title={Before you move on}]
What does \mr{कोलाहल} mean? (\textquotedblleft{}An uproar\textquotedblright{}.) Say it once more.
\end{tcolorbox}
\section[\texorpdfstring{\mr{गजबजलेला} (gajbajlelā)}{गजबजलेला (gajbajlelā)}]{\mr{गजबजलेला} (gajbajlelā) — bustling}

\label{lesson:MR-C422-gajbajlela}



\begin{quote}
Before the new one: say the Marathi for noise, then the Marathi for an uproar.
\end{quote}

\subsection*{You'll want to know: \mr{गजबजलेला}}
\textbf{\mr{गजबजलेला}} — \emph{gajbajlelā} — \textquotedblleft{}bustling\textquotedblright{}.

\begin{quote}
\mr{सणाच्या} \mr{दिवशी} \mr{बाजार} \mr{गजबजलेला} \mr{असतो}.
\end{quote}

\begin{quote}
\emph{saṇācyā divśī bājār gajbajlelā asto.}
\end{quote}

\begin{quote}
On a festival day the market is bustling.
\end{quote}

It agrees like a -\mr{ा} adjective: \mr{गजबजलेली} \mr{पेठ}, \mr{गजबजलेलं} \mr{शहर}.

\begin{cousinweb}
\textbf{\mr{गजबजलेला}} (\emph{gajbajlelā}) is the -\mr{लेला} form of the echo verb \emph{gajbajṇe}, \textquotedblleft{}to hum with people\textquotedblright{}; its origin is not settled.
\end{cousinweb}

\begin{grammarlens}[title={What you've built}]
One more word for getting around.
\end{grammarlens}

\subsection*{Your turn}
\begin{itemize}
\raggedright
  \item \emph{Say it:} \emph{gajbajlelā}
  \item \emph{Say it:} \emph{gajbajlelā}, once more
  \item \emph{Say it:} \emph{saṇācyā divśī bājār gajbajlelā asto.} — On a festival day the market is bustling.
  \item \emph{Recall:} say the Marathi for noise, then the Marathi for an uproar, then say \emph{gajbajlelā} again
\end{itemize}

\begin{tcolorbox}[breakable,colback=teal!4,colframe=teal!35!black,title={Before you move on}]
What does \mr{गजबजलेला} mean? (\textquotedblleft{}Bustling\textquotedblright{}.) Say it once more.
\end{tcolorbox}
\section[\texorpdfstring{\mr{वर्दळ} (vardaḷ)}{वर्दळ (vardaḷ)}]{\mr{वर्दळ} (vardaḷ) — bustle; coming and going}

\label{lesson:MR-C422-vardal}



\begin{quote}
Before the new one: say the Marathi for an uproar, then the Marathi for bustling.
\end{quote}

\subsection*{You'll want to know: \mr{वर्दळ}}
\textbf{\mr{वर्दळ}} — \emph{vardaḷ} — \textquotedblleft{}bustle; coming and going\textquotedblright{}.

\begin{quote}
\mr{स्टेशनबाहेर} \mr{दिवसभर} \mr{वर्दळ} \mr{असते}.
\end{quote}

\begin{quote}
\emph{sṭeśanbāher divasbhar vardaḷ aste.}
\end{quote}

\begin{quote}
Outside the station people come and go all day.
\end{quote}

Feminine.

A feminine noun: \mr{वर्दळ} \mr{कमी} \mr{झाली}, the bustle died down.

\begin{cousinweb}
\textbf{\mr{वर्दळ}} (\emph{vardaḷ}) is a Marathi word for steady coming and going; its origin is not settled.
\end{cousinweb}

\begin{grammarlens}[title={What you've built}]
One more word for getting around.
\end{grammarlens}

\subsection*{Your turn}
\begin{itemize}
\raggedright
  \item \emph{Say it:} \emph{vardaḷ}
  \item \emph{Say it:} \emph{vardaḷ}, once more
  \item \emph{Say it:} \emph{sṭeśanbāher divasbhar vardaḷ aste.} — Outside the station people come and go all day.
  \item \emph{Recall:} say the Marathi for an uproar, then the Marathi for bustling, then say \emph{vardaḷ} again
\end{itemize}

\begin{tcolorbox}[breakable,colback=teal!4,colframe=teal!35!black,title={Before you move on}]
What does \mr{वर्दळ} mean? (\textquotedblleft{}Bustle; coming and going\textquotedblright{}.) Say it once more.
\end{tcolorbox}
